% Unit 034 terjemahan Bahasa Indonesia dari chapter3.tex baris 346--621.
% Sumber: Methods of Algebra, Volume 2, commit
% 9a5803ff77dd3257484cb177f851a73770a59dd3 (CC BY 4.0).
% stable-unit-id: o014.aljabr2.chapter3.mapping-cone
% Cakupan: Bagian 3.3 lengkap.

% segment-id: o014.li2.chapter3.mapping-cone.h001
\section{Kerucut pemetaan}\label{sec:mapping-cone}
\hypertarget{o014.aljabr2.chapter3.mapping-cone}{ }
\enlargethispage{2pt}

% segment-id: o014.li2.chapter3.mapping-cone.p001
Misalkan kembali $\mathcal{A}$ suatu kategori aditif beserta kategori kompleks
yang bersesuaian, $\cate{C}(\mathcal{A})$.

% segment-id: o014.li2.chapter3.mapping-cone.d001
\begin{definisi}\label{def:Cone}
	\index{kerucut pemetaan (mapping cone)}
	\index[sym1]{Conef@$\Cone(f)$}
	Diberikan suatu morfisme $f:X\to Y$ dalam $\cate{C}(\mathcal{A})$,
	\emph{kerucut pemetaan} $\Cone(f)$ dari $f$ didefinisikan sebagai kompleks
	berikut:
% segment-id: o014.li2.chapter3.mapping-cone.q001
	\begin{align*}
		\Cone(f)^n & := X^{n+1} \oplus Y^n, \\
		d^n_{\Cone(f)} & := \;\text{notasi matriks}\;
		\begin{tikzpicture}[baseline]
			\matrix (M) [matrix of math nodes, ampersand replacement=\&, left delimiter=(, right delimiter=)] {
					-d_X^{n+1} \& 0 \\
					f^{n+1} \& d_Y^n \\
				};
				\node[above=1.2em] at (M-1-1) {\scriptsize $X^{n+1}$};
				\node[above=1.2em] at (M-1-2) {\scriptsize $Y^n$};
				\node[left=2.7em] at (M-1-1) {\scriptsize $X^{n+2}$};
				\node[left=2.7em] at (M-2-1) {\scriptsize $Y^{n+1}$};
		\end{tikzpicture} \\
		& : X^{n+1} \oplus Y^n \to X^{n+2} \oplus Y^{n+1} ,
	\end{align*}
	untuk $n\in\Z$. Karena $d_{X[1]}^n=-d_X^{n+1}$, tersedia pula notasi
	yang lebih ringkas
% segment-id: o014.li2.chapter3.mapping-cone.q002
	\[ \Cone(f) := \left( X[1] \oplus Y , \; \begin{pmatrix} d_{X[1]} & 0 \\ f[1] & d_Y \end{pmatrix} \right) . \]
\end{definisi}

% segment-id: o014.li2.chapter3.mapping-cone.p002
Mudah dibuktikan bahwa $\Cone(f)$ tetap merupakan kompleks, sebagaimana
ditunjukkan oleh perhitungan matriks berikut:
% segment-id: o014.li2.chapter3.mapping-cone.q003
\[ \begin{pmatrix} -d_X^{n+1} & 0 \\ f^{n+1} & d_Y^n \end{pmatrix} \begin{pmatrix} -d_X^n & 0 \\ f^n & d_Y^{n-1} \end{pmatrix} =
\begin{pmatrix} d_X^{n+1} d_X^n & 0 \\ - f^{n+1} d_X^n + d_Y^n f^n & d_Y^n d_Y^{n-1} \end{pmatrix}. \]

% segment-id: o014.li2.chapter3.mapping-cone.e001
\begin{contoh}
	Misalkan $f:X\to Y$ suatu morfisme dalam $\mathcal{A}$. Pandang $X,Y$
	sebagai objek $\cate{C}(\mathcal{A})$ (terkonsentrasi pada derajat $0$).
	Maka $\Cone(f)$ tepat merupakan kompleks $[X\xrightarrow{f}Y]$ (pada
	derajat $-1,0$, sedangkan semua suku lainnya sama dengan $0$).
\end{contoh}

% segment-id: o014.li2.chapter3.mapping-cone.p003
Kedua sifat funktorialitas berikut langsung diperoleh dari definisi.

% segment-id: o014.li2.chapter3.mapping-cone.pr001
\begin{proposisi}\label{prop:cone-functorial}
	Diberikan diagram komutatif dalam $\cate{C}(\mathcal{A})$
% segment-id: o014.li2.chapter3.mapping-cone.g001
	\[\begin{tikzcd}
		X \arrow[d, "f"'] \arrow[r, "\varphi"] & X' \arrow[d, "{f'}"] \\
		Y \arrow[r, "\psi"'] & Y'
	\end{tikzcd}\]
	maka $\bigl(\begin{smallmatrix} \varphi[1] & 0 \\ 0 & \psi \end{smallmatrix}\bigr)$
	memberikan morfisme $\Cone(f)\to\Cone(f')$.
\end{proposisi}

% segment-id: o014.li2.chapter3.mapping-cone.pr002
\begin{proposisi}\label{prop:Cone-A}
	Misalkan $F:\mathcal{A}\to\mathcal{A}'$ suatu funktor aditif dan
	$f:X\to Y$ suatu morfisme dalam $\cate{C}(\mathcal{A})$. Funktor yang
	bersesuaian $\cate{C}F:\cate{C}(\mathcal{A})\to
	\cate{C}(\mathcal{A}')$ memenuhi isomorfisme kanonik berikut dalam
	$\cate{C}(\mathcal{A}')$:
% segment-id: o014.li2.chapter3.mapping-cone.q004
	\[ (\cate{C}F)(\Cone(f)) \simeq \Cone(\cate{C}F(f)). \]
\end{proposisi}

% segment-id: o014.li2.chapter3.mapping-cone.p004
Kerucut pemetaan dapat ditempatkan dalam ``segitiga'' berikut, yang menjadi
landasan bagi seluruh teori selanjutnya.

% segment-id: o014.li2.chapter3.mapping-cone.d002
\begin{definisi}\label{def:cone-triangle}
	\index[sym1]{alphafbetaf@$\alpha(f), \beta(f)$}
	Tinjau suatu morfisme $f:X\to Y$ dalam $\cate{C}(\mathcal{A})$. Dari
	morfisme ini diperoleh morfisme-morfisme kanonik berikut dalam
	$\cate{C}(\mathcal{A})$:
% segment-id: o014.li2.chapter3.mapping-cone.q005
	\[ Y \xrightarrow{\alpha(f)} \Cone(f) \xrightarrow{\beta(f)} X[1] , \]
	yang secara eksplisit dinyatakan dengan matriks sebagai berikut:
% segment-id: o014.li2.chapter3.mapping-cone.q006
	\begin{align*}
		\alpha(f)^n & := \text{inklusi}\; \begin{pmatrix} 0 \\ \identity_{Y^n} \end{pmatrix}: Y^n \to X[1]^n \oplus Y^n , \\
		\beta(f)^n & := \text{proyeksi}\; \begin{pmatrix} \identity_{X[1]^n} & 0 \end{pmatrix}: X[1]^n \oplus Y^n \to X[1]^n .
	\end{align*}
	Jelas bahwa keduanya memang memberikan morfisme dalam
	$\cate{C}(\mathcal{A})$.
\end{definisi}

% segment-id: o014.li2.chapter3.mapping-cone.p005
Berikut ini dihimpun beberapa sifat homotopi kerucut pemetaan. Pertama-tama,
kita tunjukkan bahwa kerucut pemetaan dapat dipandang sebagai ``kokernel
homotopi'' atau ``kernel homotopi'' suatu morfisme. Inilah perwujudan pada
tingkat kompleks dari gagasan dasar dalam teori homotopi.

% segment-id: o014.li2.chapter3.mapping-cone.pr003
\begin{proposisi}\label{prop:homotopy-kernel-cokernel}
	Misalkan $f:X\to Y$ suatu morfisme dalam $\cate{C}(\mathcal{A})$. Terdapat
	bijeksi-bijeksi kanonik
% segment-id: o014.li2.chapter3.mapping-cone.g002
	\[\begin{tikzcd}[row sep=small]
		\Hom\left( \Cone(f), T \right) \arrow[leftrightarrow, r, "1:1"] & \left\{ (u, h): u \in \Hom(Y,T), \; h: \text{homotopi dari $uf$ ke $0$} \right\}, \\
		\Hom\left( T, \Cone(f)[-1] \right) \arrow[leftrightarrow, r, "1:1"] & \left\{ (v, k): v \in \Hom(T,X), \; k: \text{homotopi dari $fv$ ke $0$} \right\},
	\end{tikzcd}\]
	dengan $T$ sembarang objek $\cate{C}(\mathcal{A})$ dan
	$\Hom:=\Hom_{\cate{C}(\mathcal{A})}$. Secara lebih eksplisit,
% segment-id: o014.li2.chapter3.mapping-cone.l001
	\begin{compactitem}
% segment-id: o014.li2.chapter3.mapping-cone.i001
		\item citra $(u,h)$ dari $\tilde{u}:\Cone(f)\to T$ memenuhi
		$u=\tilde{u}\circ\alpha(f)$;
% segment-id: o014.li2.chapter3.mapping-cone.i002
		\item citra $(v,k)$ dari $\tilde{v}:T\to\Cone(f)[-1]$ memenuhi
		$v=\beta(f)[-1]\circ\tilde{v}$.
	\end{compactitem}
\end{proposisi}
% segment-id: o014.li2.chapter3.mapping-cone.pf001
\begin{bukti}
	Pertama-tama, tinjau kasus $\Hom\left(\Cone(f),T\right)$. Suatu morfisme
	$\tilde{u}:\Cone(f)\to T$ ekuivalen dengan keluarga morfisme dalam
	$\mathcal{A}$ berupa $h^n:X^{n+1}\to T^n$ dan $u^n:Y^n\to T^n$, dengan
	$n\in\Z$, yang memenuhi persamaan matriks
% segment-id: o014.li2.chapter3.mapping-cone.q007
	\[\begin{pmatrix}
		h^{n+1} & u^{n+1}
	\end{pmatrix} \begin{pmatrix}
		-d_X^{n+1} & 0 \\
		f^{n+1} & d_Y^n
	\end{pmatrix} = d_T^n \begin{pmatrix}
		h^n & u^n
	\end{pmatrix}. \]
	Setelah dijabarkan, persamaan ini ekuivalen dengan pernyataan bahwa
	$u=(u^n)_n:Y\to T$ merupakan morfisme dalam $\cate{C}(\mathcal{A})$ dan
	$h:=(h^{n-1})_n$ memenuhi
	$d_{\Hom^\bullet(X,T)}^{-1}h=uf$. Inilah bijeksi yang dimaksud; berdasarkan
	definisi $\alpha(f)$, persamaan $u=\tilde{u}\circ\alpha(f)$ langsung
	terlihat.

	Selanjutnya, tinjau suatu morfisme
	$\tilde{v}:T\to\Cone(f)[-1]$. Morfisme ini ekuivalen dengan keluarga
	morfisme dalam $\mathcal{A}$ berupa $v^n:T^n\to X^n$ dan
	$\underline{k}^n:T^n\to Y^{n-1}$ yang memenuhi
% segment-id: o014.li2.chapter3.mapping-cone.q008
	\[\begin{pmatrix}
		d_X^n & 0 \\
		-f^n & -d_Y^{n-1}
	\end{pmatrix} \begin{pmatrix}
		v^n \\ \underline{k}^n
	\end{pmatrix} = \begin{pmatrix}
		v^{n+1} \\ \underline{k}^{n+1}
	\end{pmatrix} d_T^n
	\quad (n \in \Z). \]
	Persamaan ini ekuivalen dengan pernyataan bahwa
	$v=(v^n)_n:T\to X$ merupakan morfisme dalam $\cate{C}(\mathcal{A})$ dan
	$k:=(-\underline{k}^n)_n$ memenuhi
	$d_{\Hom^\bullet(T,Y)}^{-1}k=fv$. Berdasarkan definisi $\beta(f)$,
	persamaan $v=\beta(f)[-1]\circ\tilde{v}$ juga langsung terlihat.
\end{bukti}

% segment-id: o014.li2.chapter3.mapping-cone.r001
\begin{catatan}[Kokernel homotopi dan kernel homotopi]\label{rem:homotopy-kernel-cokernel}
	\index{kokernel homotopi, kernel homotopi (homotopy cokernel, homotopy kernel)}
	Dari Proposisi~\ref{prop:homotopy-kernel-cokernel} diperoleh bahwa suatu
	morfisme $u:Y\to T$ (atau $v:T\to X$) memenuhi $uf$ (atau $fv$)
	homotopik nol jika dan hanya jika morfisme itu berfaktor melalui
	$\alpha(f):Y\to\Cone(f)$ (atau
	$\beta(f)[-1]:\Cone(f)[-1]\to X$). Hal ini dengan sendirinya mengingatkan
	kita pada sifat universal kokernel (atau kernel), dengan perbedaan berikut:
% segment-id: o014.li2.chapter3.mapping-cone.l002
	\begin{compactitem}
% segment-id: o014.li2.chapter3.mapping-cone.i003
		\item di sini syarat bahwa $uf$ (atau $fv$) homotopik nol menggantikan
		persamaan ketat $=0$;
% segment-id: o014.li2.chapter3.mapping-cone.i004
		\item faktorisasi $\tilde{u}$ (atau $\tilde{v}$) yang diberikan oleh
		proposisi tidak hanya membuktikan bahwa $uf$ (atau $fv$) homotopik nol,
		tetapi juga memuat cara morfisme itu dihomotopikan ke $0$; ini merupakan
		data satu tingkat lebih tinggi.
	\end{compactitem}
	Sifat-sifat ini tidak dapat ditangani begitu saja di
	$\cate{C}(\mathcal{A})$ atau $\cate{K}(\mathcal{A})$ dengan konsep-konsep
	dasar teori kategori; bahkan, kernel atau kokernel jarang ada dalam
	$\cate{K}(\mathcal{A})$. Oleh sebab itu, $Y\to\Cone(f)$ juga disebut
	\emph{kokernel homotopi} dari $f$, sedangkan
	$\beta(f)[-1]:\Cone(f)[-1]\to X$ disebut \emph{kernel homotopi} dari $f$.
	Menariknya, kerucut pemetaan $\Cone(f)$ beserta pergeserannya sekaligus
	memainkan peran ``kokernel'' dan ``kernel'' dalam pengertian ini.
\end{catatan}

% segment-id: o014.li2.chapter3.mapping-cone.p006
Suatu morfisme $f$ dalam $\cate{C}(\mathcal{A})$ disebut homotopik secara
kanonik dengan $g$ jika terdapat $h$ kanonik sedemikian sehingga
$g-f=d^{-1}h$; jika $f$ homotopik secara kanonik dengan $0$, maka morfisme
itu disebut homotopik nol secara kanonik.

% segment-id: o014.li2.chapter3.mapping-cone.pr004
\begin{proposisi}\label{prop:cone-homotopy}
	Misalkan $\mathcal{A}$ suatu kategori aditif.
% segment-id: o014.li2.chapter3.mapping-cone.l003
	\begin{enumerate}[(i)]
% segment-id: o014.li2.chapter3.mapping-cone.i005
		\item Jika $f:X\to Y$ suatu isomorfisme dalam
		$\cate{C}(\mathcal{A})$, maka $\identity_{\Cone(f)}$ homotopik nol
		secara kanonik.
% segment-id: o014.li2.chapter3.mapping-cone.i006
		\item Untuk sembarang morfisme $f:X\to Y$ dalam
		$\cate{C}(\mathcal{A})$, morfisme-morfisme dalam
		Definisi~\ref{def:cone-triangle} memenuhi
% segment-id: o014.li2.chapter3.mapping-cone.q009
		\[ \beta(f) \circ \alpha(f) = 0, \]
		sedangkan $\alpha(f)\circ f$ dan $f[1]\circ\beta(f)$ keduanya
		homotopik nol secara kanonik.
	\end{enumerate}
\end{proposisi}
% segment-id: o014.li2.chapter3.mapping-cone.pf002
\begin{bukti}
	Pertama-tama, tinjau (i). Berdasarkan funktorialitas kerucut pemetaan
	(Proposisi~\ref{prop:cone-functorial}), cukup meninjau
	$f=\identity_X$. Definisikan
% segment-id: o014.li2.chapter3.mapping-cone.q010
	\[ s^n := \begin{pmatrix} 0 & \identity_{X^n} \\ 0 & 0 \end{pmatrix} : X^{n+1} \oplus X^n \to X^n \oplus X^{n-1}, \quad n \in \Z. \]
	Perhitungan langsung memberikan
% segment-id: o014.li2.chapter3.mapping-cone.q011
	\begin{multline*}
		d^{n-1}_{\Cone(\identity_X)} s^n + s^{n+1} d^n_{\Cone(\identity_X)} = \\
		\begin{pmatrix} -d_X^n & 0 \\ \identity_{X^n} & d_X^{n-1} \end{pmatrix} \begin{pmatrix} 0 & \identity_{X^n} \\ 0 & 0 \end{pmatrix}  + \begin{pmatrix} 0 & \identity_{X^{n+1}} \\ 0 & 0 \end{pmatrix} \begin{pmatrix} -d_X^{n+1} & 0 \\ \identity_{X^{n+1}} & d_X^n \end{pmatrix}
		= \identity_{X^{n+1} \oplus X^n},
	\end{multline*}
	sehingga $s=(s^n)_{n\in\Z}$ menjadikan
	$\identity_{\Cone(\identity_X)}$ homotopik nol.

	Selanjutnya, tinjau (ii). Jelas bahwa
	$\beta(f)\circ\alpha(f)=0$. Homotopi-homotopi yang tersisa diperoleh dari
	Proposisi~\ref{prop:homotopy-kernel-cokernel}: ambil $T=\Cone(f)$, maka
	$\tilde{u}:=\identity_{\Cone(f)}\in\Hom(\Cone(f),T)$ menjadikan
	$\alpha(f)\circ f$ homotopik nol; ambil $T=\Cone(f)[-1]$, maka
	$\tilde{v}:=\identity_{\Cone(f)[-1]}\in\Hom(T,\Cone(f)[-1])$ menjadikan
	$f\circ\beta(f)[-1]$ homotopik nol, dan dengan demikian juga menjadikan
	$f[1]\circ\beta(f)$ homotopik nol.
\end{bukti}

% segment-id: o014.li2.chapter3.mapping-cone.p007
Definisi~\ref{def:cone-triangle} menghasilkan dua morfisme baru
$\alpha(f)$ dan $\beta(f)$ dari $f$. Hingga homotopi, kerucut pemetaan kedua
morfisme itu tidak menghasilkan objek baru; kita mulai dengan
$\Cone(\alpha(f))$ untuk menjelaskan hal tersebut. Pertama-tama, perhatikan
bahwa
% segment-id: o014.li2.chapter3.mapping-cone.q012
\[ \Cone(\alpha(f))^n = Y^{n+1} \oplus \Cone(f)^n = Y^{n+1} \oplus X^{n+1} \oplus Y^n. \]
Dalam notasi matriks,
% segment-id: o014.li2.chapter3.mapping-cone.q013
\begin{align*}
	\alpha(\alpha(f)) & = \begin{pmatrix}
		0 & 0 \\
		\identity_{X[1]} & 0 \\
		0 & \identity_Y
	\end{pmatrix} : \Cone(f) \to \Cone(\alpha(f)) , \\
	\beta(\alpha(f)) & = \begin{pmatrix}
		\identity_{Y[1]} & 0 & 0
	\end{pmatrix} : \Cone(\alpha(f)) \to Y[1].
\end{align*}

% segment-id: o014.li2.chapter3.mapping-cone.le001
\begin{lema}\label{prop:cone-alpha}
	Untuk suatu morfisme $f:X\to Y$ dalam $\cate{C}(\mathcal{A})$, notasi
	matriks mendefinisikan sepasang morfisme
% segment-id: o014.li2.chapter3.mapping-cone.g003
	\[ \begin{pmatrix} -f[1] \\ \identity_{X[1]} \\ 0 \end{pmatrix} : \begin{tikzcd} X[1] \arrow[r, shift left, "\phi"] & \Cone(\alpha(f)) \arrow[l, shift left, "\psi"] \end{tikzcd} : \begin{pmatrix} 0 & \identity_{X[1]} & 0 \end{pmatrix}. \]
	Keduanya memenuhi $\psi\circ\phi=\identity_{X[1]}$ dan menjadikan diagram
	berikut komutatif dalam $\cate{K}(\mathcal{A})$:
% segment-id: o014.li2.chapter3.mapping-cone.g004
	\[\begin{tikzcd}
		\Cone(f) \arrow[d, "{\identity_{\Cone(f)}}"'] \arrow[r, "\beta(f)"] & X[1] \arrow[d, "\phi"] \arrow[r, "{-f[1]}"] & Y[1] \arrow[d, "{\identity_{Y[1]}}"] \\
		\Cone(f) \arrow[r, "{\alpha(\alpha(f))}"'] & \Cone(\alpha(f)) \arrow[r, "\beta(\alpha(f))"'] & Y[1]
	\end{tikzcd}\]
	Selain itu, $\phi,\psi$ saling invers dalam $\cate{K}(\mathcal{A})$.
\end{lema}
% segment-id: o014.li2.chapter3.mapping-cone.pf003
\begin{bukti}
	Berdasarkan uraian sebelumnya tentang $\Cone(\alpha(f))$, dalam notasi
	matriks $d_{\Cone(\alpha(f))}^n$ diberikan oleh
% segment-id: o014.li2.chapter3.mapping-cone.q014
	\[ \begin{pmatrix} -d_Y^{n+1} & 0 & 0 \\ 0 & -d_X^{n+1} & 0 \\ \identity_{Y^{n+1}} & f^{n+1} & d_Y^n \end{pmatrix} : Y^{n+1} \oplus X^{n+1} \oplus Y^n \to Y^{n+2} \oplus X^{n+2} \oplus Y^{n+1}. \]
	Cukup diverifikasi pernyataan-pernyataan berikut dalam
	$\cate{C}(\mathcal{A})$:
% segment-id: o014.li2.chapter3.mapping-cone.l004
	\begin{itemize}
% segment-id: o014.li2.chapter3.mapping-cone.i007
		\item $\phi:=(\phi^n)_{n\in\Z}$ dan $\psi:=(\psi^n)_{n\in\Z}$
		keduanya merupakan morfisme kompleks;
% segment-id: o014.li2.chapter3.mapping-cone.i008
		\item $\psi\circ\phi=\identity_{X[1]}$;
% segment-id: o014.li2.chapter3.mapping-cone.i009
		\item $\psi\circ\alpha(\alpha(f))=\beta(f)$;
% segment-id: o014.li2.chapter3.mapping-cone.i010
		\item $\beta(\alpha(f))\circ\phi=-f[1]$;
% segment-id: o014.li2.chapter3.mapping-cone.i011
		\item terdapat
		$s=(s^n)_{n\in\Z}\in\Hom^{-1}\left(\Cone(\alpha(f)),
		\Cone(\alpha(f))\right)$ sedemikian sehingga
		$\identity_{\Cone(\alpha(f))}-\phi\circ\psi
		=d^{-1}_{\Hom^\bullet}(s)$.
	\end{itemize}
	Empat pernyataan pertama bersifat elementer. Untuk pernyataan terakhir,
	ambillah, dalam notasi matriks,
% segment-id: o014.li2.chapter3.mapping-cone.q015
	\[ s^n := \begin{pmatrix} 0 & 0 & \identity_{Y^n} \\ 0 & 0 & 0 \\ 0 & 0 & 0 \end{pmatrix}: \Cone(\alpha(f))^n \to \Cone(\alpha(f))^{n-1} \]
	dan lakukan verifikasi langsung.
\end{bukti}

% segment-id: o014.li2.chapter3.mapping-cone.p008
Untuk kasus $\beta(f)$, kita melakukan sedikit modifikasi dan meninjau
kerucut pemetaan dari $-\beta(f)[-1]:\Cone(f)[-1]\to X$.

% segment-id: o014.li2.chapter3.mapping-cone.d003
\begin{definisi}\label{def:Cyl}
	\index{silinder pemetaan (mapping cylinder)}
	\index[sym1]{Cylf@$\Cyl(f)$}
	Untuk suatu morfisme $f:X\to Y$ dalam $\cate{C}(\mathcal{A})$,
	\emph{silinder pemetaan} dari $f$ didefinisikan sebagai
% segment-id: o014.li2.chapter3.mapping-cone.q016
	\[ \Cyl(f) := \Cone\left( \Cone(f)[-1] \xrightarrow{-\beta(f)[-1]} X \right). \]
	Silinder ini dilengkapi morfisme-morfisme kanonik
	$X\to\Cyl(f)\to\Cone(f)$.
\end{definisi}

% segment-id: o014.li2.chapter3.mapping-cone.p009
Silinder pemetaan mempunyai funktorialitas yang bentuknya sama dengan
funktorialitas kerucut pemetaan (Proposisi~\ref{prop:cone-functorial}).
Berdasarkan Catatan~\ref{rem:homotopy-kernel-cokernel}, $\Cyl(f)$ dapat
dibandingkan dengan kocitra dalam kategori aditif---objek ini kurang lebih
merupakan ``kokernel dari kernel'' (Proposisi~\ref{prop:Im-Coim-additive})---
sehingga dapat pula dipandang sebagai kocitra homotopi dari $f$.

% segment-id: o014.li2.chapter3.mapping-cone.p010
Secara lebih eksplisit,
$\Cyl(f)^n=\Cone(f)^n\oplus X^n=X^{n+1}\oplus Y^n\oplus X^n$; morfisme
$X\to\mathrm{Cyl}(f)$ (atau $\mathrm{Cyl}(f)\to\Cone(f)$) ialah inklusi ke
suku ketiga pada jumlah langsung (atau proyeksi ke dua suku pertama), dan
% segment-id: o014.li2.chapter3.mapping-cone.q017
\[ d_{\Cyl(f)}^n =
	\begin{pmatrix} d_{\Cone(f)}^n & 0 \\ -\beta(f)^n & d_X^n \end{pmatrix}
	= \begin{pmatrix}
		-d_X^{n+1} & 0 & 0 \\
		f^{n+1} & d_Y^n & 0 \\
		-\identity_{X^{n+1}} & 0 & d_X^n
\end{pmatrix} .\]

% segment-id: o014.li2.chapter3.mapping-cone.le002
\begin{lema}\label{prop:cone-beta}
	Untuk suatu morfisme $f:X\to Y$ dalam $\cate{C}(\mathcal{A})$, dapat
	didefinisikan sepasang morfisme
% segment-id: o014.li2.chapter3.mapping-cone.g005
	\[ \begin{pmatrix} 0 \\ \identity_Y \\ 0 \end{pmatrix} : \begin{tikzcd} Y \arrow[r, shift left, "\phi"] & \Cyl(f) \arrow[l, shift left, "\psi"] \end{tikzcd} : \begin{pmatrix} 0 & \identity_Y & f \end{pmatrix}. \]
	Keduanya memenuhi $\psi\circ\phi=\identity_Y$ dan menjadikan diagram
	berikut komutatif dalam $\cate{K}(\mathcal{A})$:
% segment-id: o014.li2.chapter3.mapping-cone.g006
	\[\begin{tikzcd}
		X \arrow[d, "{\identity_X}"'] \arrow[r, "f"] & Y \arrow[d, "\phi"] \arrow[r, "{\alpha(f)}"] & \Cone(f) \arrow[d, "{\identity_{\Cone(f)}}"] \\
		X \arrow[r] & \Cyl(f) \arrow[r] & \Cone(f)
	\end{tikzcd}\]
	Selain itu, $\phi,\psi$ saling invers dalam $\cate{K}(\mathcal{A})$.
Lebih lanjut, $f$ berfaktor sebagai komposisi
$X\to\Cyl(f)\xrightarrow{\psi}Y$.
\end{lema}
% segment-id: o014.li2.chapter3.mapping-cone.pf004
\begin{bukti}
	Pembaca dapat memverifikasi langsung dari definisi bahwa $\phi$ dan $\psi$
	memang memberikan morfisme kompleks. Setelah itu,
	$\psi\circ\phi=\identity_Y$ langsung terlihat. Selanjutnya, untuk
	membuktikan bahwa $\phi$ dan $\psi$ saling invers dalam
	$\cate{K}(\mathcal{A})$, seperti pada Lema~\ref{prop:cone-alpha}, ambillah
% segment-id: o014.li2.chapter3.mapping-cone.q018
	\begin{align*}
		s^n & := \begin{pmatrix} 0 & 0 & -\identity_{X^n} \\ 0 & 0 & 0 \\ 0 & 0 & 0 \end{pmatrix}: X^{n+1} \oplus Y^n \oplus X^n \to X^n \oplus Y^{n-1} \oplus X^{n-1}, \\
		s & := (s^n)_{n \in \Z},
	\end{align*}
	dan verifikasi
	$\identity_{\Cyl(f)}-\phi\circ\psi=d^{-1}_{\Hom^\bullet}(s)$.

	Kotak kanan pada diagram sudah komutatif dalam $\cate{C}(\mathcal{A})$.
	Untuk kotak kiri, cukup perhatikan bahwa komposisi
	$X\to\mathrm{Cyl}(f)\xrightarrow{\psi}Y$ sama dengan $f$.
\end{bukti}

% segment-id: o014.li2.chapter3.mapping-cone.p011
Lema~\ref{prop:cone-alpha} bersama Lema~\ref{prop:cone-beta} menunjukkan
bahwa, hingga isomorfisme dalam $\cate{K}(\mathcal{A})$, sembarang morfisme
$f:X\to Y$ dapat diganti dengan proyeksi yang jauh lebih sederhana
$\Cone(\alpha(f))[-1]\to Y$ (atau inklusi $X\to\Cyl(f)$), dan konstruksi ini
funktorial dalam $f$. Inilah suatu gagasan penting dalam aljabar homologis
dan teori homotopi.

% segment-id: o014.li2.chapter3.mapping-cone.p012
Kasus khusus silinder pemetaan dengan $f=\identity_X$ mempunyai kegunaan lain:
silinder itu dapat digunakan untuk menafsirkan homotopi. Bagi pembaca yang
mengenal teori homologi, semua ini mempunyai tafsiran topologis, tetapi di
sini kita hanya membahas versi aljabarnya.

% segment-id: o014.li2.chapter3.mapping-cone.pr005
\begin{proposisi}\label{prop:cylinder-homotopy}
	\index[sym1]{CylX@$\Cyl_X$}
	Misalkan $X$ suatu objek $\cate{C}(\mathcal{A})$ dan definisikan
	$\mathrm{Cyl}_X:=\mathrm{Cyl}(\identity_X)$. Perhatikan bahwa
	$\mathrm{Cyl}_X^n=X^{n+1}\oplus X^n\oplus X^n$.
% segment-id: o014.li2.chapter3.mapping-cone.l005
	\begin{enumerate}[(i)]
% segment-id: o014.li2.chapter3.mapping-cone.i012
		\item Dalam $\cate{C}(\mathcal{A})$ terdapat morfisme-morfisme
% segment-id: o014.li2.chapter3.mapping-cone.g007
		$\begin{tikzcd}
			X \arrow[shift left, r, "i_0"] \arrow[shift right, r, "i_1"'] & \Cyl_X \arrow[r, "j"] & X
		\end{tikzcd}$,
		yang dalam notasi matriks didefinisikan untuk setiap $n\in\Z$ oleh
% segment-id: o014.li2.chapter3.mapping-cone.q019
		\begin{equation*}
			i_0^n := \begin{pmatrix} 0 \\ 0 \\ \identity_{X^n} \end{pmatrix} , \quad
			i_1^n := \begin{pmatrix} 0 \\ \identity_{X^n} \\ 0 \end{pmatrix} , \quad j^n := \begin{pmatrix} 0 & \identity_{X^n} & \identity_{X^n} \end{pmatrix},
		\end{equation*}
		Morfisme-morfisme ini memenuhi $ji_0=\identity_X=ji_1$, dan citra $j$
		dalam $\cate{K}(\mathcal{A})$ merupakan isomorfisme.
% segment-id: o014.li2.chapter3.mapping-cone.i013
		\item Untuk sembarang objek $Y$ dari $\cate{C}(\mathcal{A})$, terdapat
		bijeksi
% segment-id: o014.li2.chapter3.mapping-cone.q020
		\begin{align*}
			\left\{\begin{array}{r|l}
				(f, g, h) & f, g \in \Hom_{\cate{C}(\mathcal{A})}(X, Y) \\
				& h \in \Hom^{-1}(X, Y) \\
				& g - f = d^{-1}_{\Hom^\bullet(X, Y)} h
			\end{array}\right\} & \xrightarrow{1:1} \Hom_{\cate{C}(\mathcal{A})}\left( \mathrm{Cyl}_X, Y \right) \\
			(f, g, h) & \longmapsto \tilde{h} = (\tilde{h}^n)_{n \in \Z}, \; \tilde{h}^n := \begin{pmatrix} h^n & g^n & f^n \end{pmatrix}.
		\end{align*}
		Secara khusus, $f=\tilde{h}i_0$ dan $g=\tilde{h}i_1$.
	\end{enumerate}
\end{proposisi}
% segment-id: o014.li2.chapter3.mapping-cone.pf005
\begin{bukti}
	Untuk (i), perhatikan bahwa $i_0$ tidak lain ialah morfisme kanonik
	$X\to\mathrm{Cyl}(\identity_X)$ dalam Definisi~\ref{def:Cyl}, sedangkan
	$i_1$ ialah morfisme $\phi:X\to\mathrm{Cyl}(\identity_X)$ dalam
	Lema~\ref{prop:cone-beta}. Diagram komutatif dalam
	Lema~\ref{prop:cone-beta} menunjukkan bahwa keduanya memberikan
	isomorfisme yang sama dalam $\cate{K}(\mathcal{A})$.

	Untuk $j$, mudah diverifikasi bahwa
	$j^{n+1}d_{\Cyl_X}^n=\bigl(0\;d_X^n\;d_X^n\bigr)=d_X^nj^n$.
	Jadi, $j$ memang merupakan morfisme, sedangkan
	$ji_0=\identity_X=ji_1$ langsung terlihat. Karena $i_0$ dan $i_1$
	merupakan isomorfisme dalam $\cate{K}(\mathcal{A})$, demikian pula $j$.

	Untuk (ii), tuliskan sembarang unsur $\tilde{h}$ dari
	$\Hom^0(\mathrm{Cyl}_X,Y)$ dalam bentuk
	\[
		\tilde{h}=\left((h^n,g^n,f^n)\right)_{n\in\Z},
	\]
	dengan
% segment-id: o014.li2.chapter3.mapping-cone.q021
	\[ h^n: X^{n+1} \to Y^n, \quad f^n: X^n \to Y^n, \quad g^n: X^n \to Y^n \]
	semuanya merupakan morfisme dalam $\mathcal{A}$. Dengan menghilangkan
	superskrip dan menggunakan notasi matriks, hitung
% segment-id: o014.li2.chapter3.mapping-cone.q022
	\begin{align*}
		\tilde{h} \; d_{\mathrm{Cyl}_X} & = \begin{pmatrix} h & g & f \end{pmatrix} \begin{pmatrix} -d_X & 0 & 0 \\ \identity_X & d_X & 0 \\ -\identity_X & 0 & d_X \end{pmatrix} = \begin{pmatrix} -hd_X + g - f & g d_X & f d_X \end{pmatrix} , \\
		d_Y \; \tilde{h} & = \begin{pmatrix} d_Y h & d_Y g & d_Y f \end{pmatrix}.
	\end{align*}
	Oleh karena itu, $\tilde{h}$ merupakan morfisme kompleks jika dan hanya
	jika $f,g\in\Hom_{\cate{C}(\mathcal{A})}(X,Y)$ dan
	$g-f=hd_X+d_Yh$.
\end{bukti}

% segment-id: o014.li2.chapter3.mapping-cone.r002
\begin{catatan}\label{rem:cone-chain}
	\index{kerucut pemetaan}\index{silinder pemetaan}
	Kerucut pemetaan dan silinder pemetaan tentu saja mempunyai versi untuk
	kompleks rantai (lihat Catatan~\ref{rem:cochain-vs-chain}). Diberikan suatu
	morfisme kompleks rantai $f:X\to Y$, definisikan
% segment-id: o014.li2.chapter3.mapping-cone.q023
	\[\begin{array}{rlrl}
		\Cone(f)_n & :=(X[-1] \oplus Y)_n, &
		\Cyl(f)_n & := (X[-1] \oplus Y \oplus X)_n, \\
		d_{\Cone(f)} & := \begin{pmatrix} d_{X[-1]} & 0 \\ f[-1] & d_Y \end{pmatrix}, &
		d_{\Cyl(f)} & := \begin{pmatrix} d_{X[-1]} & 0 & 0 \\ f[-1] & d_Y & 0 \\ -\identity_{X[-1]} & 0 & d_X \end{pmatrix}.
	\end{array}\]

	Semua pernyataan dalam bagian ini dapat dipindahkan ke kasus kompleks
	rantai. Secara khusus, terdapat morfisme-morfisme kanonik
% segment-id: o014.li2.chapter3.mapping-cone.q024
	\[ X \xrightarrow{f} Y \xrightarrow{\alpha(f)} \Cone(f) \xrightarrow{\beta(f)} X[-1], \quad X \to \Cyl(f) \to \Cone(f). \]
\end{catatan}
